% Instructor-authored original material. February 2024 LSAT/LSAC identifiers
% are used only to identify skill inspiration. This material is not endorsed by LSAC.
%
% SOURCE MAP (official form -> general skill only)
% 2024 Feb RC Q1  -> primary purpose
% 2024 Feb RC Q2  -> passage-supported inference
% 2024 Feb RC Q3  -> stated detail
% 2024 Feb RC Q4  -> author's evaluative emphasis
% 2024 Feb RC Q5  -> EXCEPT / detail inventory
% 2024 Feb RC Q6  -> viewpoint inference
% 2024 Feb RC Q7  -> shared primary concern across paired passages
% 2024 Feb RC Q8  -> information in Passage B but not Passage A
% 2024 Feb RC Q9  -> function of a description in Passage A
% 2024 Feb RC Q10 -> weaken an explanation in Passage B
% 2024 Feb RC Q11 -> statement supported by both passages
% 2024 Feb RC Q12 -> likely agreement between authors
% 2024 Feb RC Q13 -> main point of Passage A
% 2024 Feb RC Q14 -> analogy for the relationship between approaches
% 2024 Feb RC Q15 -> central idea
% 2024 Feb RC Q16 -> stated detail
% 2024 Feb RC Q17 -> paragraph function
% 2024 Feb RC Q18 -> sentence function
% 2024 Feb RC Q19 -> analogy for a theory's status
% 2024 Feb RC Q20 -> author's likely agreement
% 2024 Feb RC Q21 -> strengthen a theory
% 2024 Feb RC Q22 -> main point
% 2024 Feb RC Q23 -> author's criticism of a theory
% 2024 Feb RC Q24 -> implication of a theory
% 2024 Feb RC Q25 -> purpose of a conceptual distinction
% 2024 Feb RC Q26 -> meaning in context
% 2024 Feb RC Q27 -> author's evaluation of a theory

\begingroup
\renewcommand{\labelenumii}{(\Alph{enumii})}

\section*{Reading Comprehension: Original Analogues}

\begin{center}
\fbox{\parbox{0.92\linewidth}{\textbf{Instructor-authored original material.}
The LSAT and LSAC names appear only to identify the source of general skill inspiration.
The passages, facts, questions, and answer choices are new. They are not LSAC material
and are not endorsed by LSAC.}}
\end{center}

Each packet is designed for approximately 30--50 minutes of collaborative reading,
individual selection, and group justification.

\subsection*{Packet 1: Restoring Alder Creek (Questions 1--6)}

For much of the twentieth century, Alder Creek crossed the city through a straight
concrete channel. The channel moved storm water rapidly, but it also accelerated water
downstream, provided little habitat, and separated nearby soil from the creek. When civil
engineer Marisol Vega proposed replacing parts of it with what she called a ``slow-water
corridor,'' critics imagined an attempt to reconstruct a vanished wilderness in the middle
of a dense neighborhood. Vega's proposal was less nostalgic. She treated the creek as a
working urban system whose ecological and drainage functions could be improved without
pretending that streets and buildings would disappear.

The first pilot reach used low woven-willow hurdles to redirect ordinary flows, gravel
riffles to introduce oxygen, and shallow benches that could temporarily hold storm water.
The hurdles were deliberately permeable. During a major storm, water could pass through
and over them rather than accumulating behind a rigid wall. Sensors installed above and
below the pilot showed that peak flow downstream fell by 18 percent during moderate storms.
Moisture remained in adjacent soil several days longer, and surveys recorded the return of
three native aquatic-insect species. Fish populations did not increase during the two-year
study, however, and two road culverts still interrupted movement along the creek.

Vega regarded those mixed results as information rather than as a verdict. The project had
never promised to restore a single historical condition. In a city, rainfall patterns,
pavement, shade, and public use keep changing. The design therefore included monthly
inspection by neighborhood crews. After observing bank erosion, a crew could reposition a
hurdle; after a dry season, it could alter a shallow bench. Sensor readings and residents'
observations were discussed together at open meetings before any substantial adjustment.

This maintenance requirement is sometimes presented as a weakness. A concrete channel,
after all, can appear finished. But its apparent permanence merely hides decisions about
where water and risk are sent. Alder Creek's visible need for care makes those decisions
available for public revision. The pilot's greatest achievement is consequently not any
single percentage or species count. It is the combination of measurable ecological gains
with a local capacity to notice failure and change course. Other cities should not copy
every hurdle or bench. They should copy the commitment to treat restoration as an
evidence-guided practice rather than a completed object.

\begin{enumerate}
% Source map: 2024 Feb RC Q1 -- primary purpose.
\item The passage is primarily concerned with
\begin{enumerate}
\item proving that concrete channels always increase urban flooding
\item describing the history of every redesign of Alder Creek
\item explaining why an adaptive restoration project offers a useful model despite incomplete results
\item arguing that urban streams should be returned to one exact historical condition
\item comparing the cost of willow hurdles with the cost of road culverts
\end{enumerate}

% Source map: 2024 Feb RC Q2 -- inference.
\item Which statement is most strongly supported by the passage?
\begin{enumerate}
\item Vega would judge a project successful only if fish populations increased immediately.
\item A design can count as restoration even when it does not recreate a pre-urban landscape.
\item Neighborhood observations were collected but excluded from design decisions.
\item Rigid barriers are always safer than permeable structures during major storms.
\item The pilot eliminated every interruption to movement along the creek.
\end{enumerate}

% Source map: 2024 Feb RC Q3 -- stated detail.
\item According to the passage, what was one intended function of the gravel riffles?
\begin{enumerate}
\item introducing oxygen into the water
\item preventing residents from reaching the creek
\item replacing the road culverts
\item measuring downstream peak flow
\item supporting the willow hurdles during storms
\end{enumerate}

% Source map: 2024 Feb RC Q4 -- evaluative emphasis.
\item Which feature of the project does the author value most highly?
\begin{enumerate}
\item the appearance of a finished, permanent design
\item the return of fish within the study period
\item the removal of all pavement near the creek
\item the union of measurable improvement with community capacity to revise the design
\item the use of a single standard design that other cities can copy exactly
\end{enumerate}

% Source map: 2024 Feb RC Q5 -- EXCEPT.
\item The passage mentions each of the following as part of the project or its monitoring \emph{except}
\begin{enumerate}
\item woven-willow hurdles
\item gravel riffles
\item shallow floodwater benches
\item downstream sensors
\item stocking the creek with hatchery fish
\end{enumerate}

% Source map: 2024 Feb RC Q6 -- viewpoint inference.
\item Vega would most likely agree with which statement?
\begin{enumerate}
\item A restoration design should remain unchanged once construction ends.
\item Failure to reproduce a historical ecosystem makes an urban project meaningless.
\item A project's need for continuing revision can be a source of strength.
\item Resident observations should never be considered alongside sensor data.
\item Ecological goals should take priority over drainage under every condition.
\end{enumerate}
\end{enumerate}

\subsection*{Packet 2: Citizen Science Infrastructure (Questions 7--14)}

\textbf{Passage A.}
The Open Air Shelf began as a small lending counter for portable air-quality sensors.
Its founders assumed that access to instruments was the main obstacle preventing residents
from documenting pollution. The first year appeared to confirm their view: hundreds of
measurements were uploaded from neighborhoods beyond the reach of official monitoring
stations. Yet many files lacked calibration records, and participants often could not tell
whether a short spike reflected traffic, weather, or a faulty sensor.

The program responded by becoming more than what one coordinator called ``a hardware
library.'' Borrowers now attend a brief calibration clinic, attach field notes to readings,
and use a shared dashboard that displays wind direction and nearby official measurements.
Volunteer reviewers flag unusual results for follow-up rather than silently deleting them.
The original lending model remains important because individually purchased instruments
would be too costly for many residents. But access alone was only a first step. By adding
interpretation and review, the program turned isolated measurements into evidence that
neighborhood groups and city scientists can evaluate together.

\textbf{Passage B.}
Tide Cart serves fishing towns that are too small to maintain permanent marine laboratories.
Its van visits each harbor every week with salinity probes, microscopes, and a satellite
connection. Residents may bring water samples, borrow instruments until the next visit,
or compare their observations with an immediately updated regional map. Local stewards,
trained by the program, check instruments between visits and help newcomers record sampling
conditions consistently.

Tide Cart's organizers attribute its unusually high volunteer retention to rapid feedback.
Participants see their measurements on the regional map while the event they are observing
is still unfolding, so they can return to a site if a result looks surprising. The service
keeps costs low by using one mobile laboratory for several towns and by recruiting stewards
locally. Harbor associations pay modest subscriptions, which organizers expect will cover
routine operations once the program reaches twelve towns. Tide Cart is valuable not merely
because it circulates equipment, but because it makes careful interpretation available where
a fixed laboratory would be impractical.

\begin{enumerate}\setcounter{enumi}{6}
% Source map: 2024 Feb RC Q7 -- shared primary concern.
\item Both passages are primarily concerned with
\begin{enumerate}
\item replacing professional scientists with unpaid residents
\item improving citizen science by pairing instruments with interpretive support
\item proving that mobile laboratories always cost less than fixed ones
\item opposing the use of official environmental measurements
\item showing that calibration is unnecessary for local observations
\end{enumerate}

% Source map: 2024 Feb RC Q8 -- B but not A.
\item Which feature is mentioned in Passage B but not in Passage A?
\begin{enumerate}
\item lending scientific instruments
\item training participants to record observations consistently
\item comparing local observations with a larger data set
\item transporting a laboratory among several communities on a regular route
\item helping participants interpret surprising measurements
\end{enumerate}

% Source map: 2024 Feb RC Q9 -- function of description.
\item Passage A calls the early program ``a hardware library'' chiefly to
\begin{enumerate}
\item contrast its original limited service with its later, broader support
\item criticize libraries for lending objects other than books
\item show that the founders intended to exclude city scientists
\item explain why the sensors were too expensive to purchase
\item suggest that the program should stop lending instruments
\end{enumerate}

% Source map: 2024 Feb RC Q10 -- weaken an explanation.
\item Which fact would most weaken Passage B's explanation for Tide Cart's high volunteer retention?
\begin{enumerate}
\item Satellite service is occasionally interrupted by storms.
\item Most Tide Cart volunteers belonged to long-standing conservation groups before joining, whereas comparison programs recruited first-time volunteers.
\item Some towns collect more samples than other towns.
\item The regional map contains official as well as volunteer measurements.
\item Local stewards receive additional instrument training.
\end{enumerate}

% Source map: 2024 Feb RC Q11 -- supported by both.
\item Which statement is supported by both passages?
\begin{enumerate}
\item Useful community data can require support beyond mere access to equipment.
\item Citizen-science programs should avoid all government measurements.
\item Subscription fees are the only sustainable funding method.
\item Mobile laboratories produce more accurate data than fixed laboratories.
\item Every unusual measurement is caused by faulty equipment.
\end{enumerate}

% Source map: 2024 Feb RC Q12 -- author agreement.
\item The authors would most likely agree that
\begin{enumerate}
\item residents cannot learn to calibrate scientific instruments
\item equipment lending has value only in large cities
\item immediate publication guarantees that a measurement is accurate
\item instruments are most useful when participants can contextualize what they record
\item local volunteers should make environmental policy without professional input
\end{enumerate}

% Source map: 2024 Feb RC Q13 -- main point of A.
\item Which choice best states the main point of Passage A?
\begin{enumerate}
\item Open Air Shelf failed because residents could not afford personal sensors.
\item Instrument lending produced no useful measurements until city scientists took control.
\item An initially valuable lending program became more useful by adding calibration, context, and review.
\item The shared dashboard should replace field notes and official monitoring stations.
\item Volunteer reviewers should delete every unusual result before publication.
\end{enumerate}

% Source map: 2024 Feb RC Q14 -- relationship analogy.
\item The relationship between Open Air Shelf's original model and the later Open Air Shelf and Tide Cart models is most like the relationship between
\begin{enumerate}
\item a discarded prototype and an unrelated invention that happens to share its name
\item a neighborhood store and a chain that sells entirely different goods
\item a useful map collection and newer map services that retain the maps while adding navigation and interpretation
\item a teacher and students who repeat every lesson without modification
\item an expensive product and cheaper imitations that remove its central feature
\end{enumerate}
\end{enumerate}

\subsection*{Packet 3: The Black Bands of Kestrel Cavern (Questions 15--21)}

Thin black bands recur inside the pale mineral layers of Kestrel Cavern. Because each layer
formed as mineral-rich water dripped from the ceiling, researchers first treated the bands
as ordinary traces of sediment carried by seasonal floods. Microscopy later showed that the
dark material contains unusually high concentrations of carbon and almost no river clay.
The finding prompted two different explanations, neither of which fits comfortably within
the usual account of how cave deposits grow.

The first explanation attributes the bands to smoke from ancient hearths. On this view,
fine soot entered cracks above the occupied chamber and became trapped as new mineral formed.
Radiocarbon dates from several bands roughly overlap periods when artifacts indicate human
occupation. The account also explains why the dark layers appear intermittently. A mass-balance
calculation presents a difficulty, however: the hearths currently documented in the cavern
would have produced only about one tenth of the carbon measured in the bands. Supporters
suggest that undiscovered fires or unusually efficient transport through cracks may close
the gap, but neither proposal has yet been confirmed.

The second explanation invokes microorganisms that live within slowly growing mineral.
Related microbes collected from deep cracks produce a dark carbon-rich film when oxygen is
scarce. If their ancient relatives behaved similarly, periodic changes in groundwater could
have triggered growth of films that were later sealed into the deposit. This hypothesis
would require researchers to revise the familiar picture of cave formations as passive
chemical records: at least some layers would partly record biological activity. No intact
cells have been recovered from the old bands, so the mechanism remains inferential.

The microbial proposal is nevertheless unusually testable for such a revisionary idea.
Laboratory cultures produce a distinctive ratio of two stable carbon isotopes. Researchers
can search for the same ratio in Kestrel's bands and in young deposits from nearby unoccupied
caves. A match would not by itself eliminate every possible combustion source, but it would
give the microbial account a prediction that can be checked without waiting for another
ancient layer to form. In that respect, the proposal's radical implications may be offset
by its local accessibility to experiment.

\begin{enumerate}\setcounter{enumi}{14}
% Source map: 2024 Feb RC Q15 -- central idea.
\item Which choice best expresses the central idea of the passage?
\begin{enumerate}
\item Seasonal flooding has been conclusively established as the source of Kestrel's black bands.
\item Two explanations for the bands have empirical attractions and unresolved problems.
\item Microorganisms form every mineral layer found in caves.
\item Ancient people never occupied Kestrel Cavern.
\item Isotope ratios cannot distinguish biological carbon from combustion carbon.
\end{enumerate}

% Source map: 2024 Feb RC Q16 -- stated detail.
\item The passage states that related living microbes
\begin{enumerate}
\item consume river clay during seasonal floods
\item produce a dark film under low-oxygen conditions
\item survive only in occupied caves
\item generate ten times more carbon than a hearth
\item have been recovered intact from the ancient bands
\end{enumerate}

% Source map: 2024 Feb RC Q17 -- paragraph function.
\item The second paragraph primarily serves to
\begin{enumerate}
\item describe the combustion explanation and identify a quantitative difficulty for it
\item prove that human occupation began before the cavern formed
\item show that radiocarbon dating is unreliable in caves
\item reject microscopy as a method for studying mineral layers
\item explain how modern microbes produce isotope ratios
\end{enumerate}

% Source map: 2024 Feb RC Q18 -- sentence function.
\item The final sentence chiefly
\begin{enumerate}
\item reports that experiments have already confirmed the microbial account
\item identifies testability as a consideration favoring further study of a radical proposal
\item concedes that the combustion explanation cannot be tested
\item states that nearby caves have identical black bands
\item summarizes the history of human occupation at Kestrel
\end{enumerate}

% Source map: 2024 Feb RC Q19 -- analogous theory status.
\item Which situation is most analogous to the status of the microbial explanation?
\begin{enumerate}
\item A linguistic theory preserves every accepted assumption but has no observable consequences.
\item A medical hypothesis repeats an established explanation using new terminology.
\item A historical claim conflicts with dated records and offers no way to resolve the conflict.
\item A new account of memory revises a basic model but predicts a pattern measurable in current laboratory tasks.
\item A geological proposal explains an observation only by declaring the observation mistaken.
\end{enumerate}

% Source map: 2024 Feb RC Q20 -- author agreement.
\item The author would most likely agree that
\begin{enumerate}
\item nearby young cave deposits could provide evidence relevant to an ancient process
\item a testable theory must be correct
\item no possible fire could produce the isotope ratio in question
\item the missing carbon proves that no hearths existed
\item radical implications count against a theory regardless of its evidence
\end{enumerate}

% Source map: 2024 Feb RC Q21 -- strengthen a theory.
\item Which finding would most strengthen the combustion explanation?
\begin{enumerate}
\item The bands contain minerals found throughout the cavern.
\item Carbon particles in the bands match experimentally burned local wood, and the densest bands align with independently dated occupation peaks.
\item Living microbes occur in several deep cracks.
\item Floodwater sometimes enters the lowest chamber.
\item Pale layers also contain small amounts of carbon.
\end{enumerate}
\end{enumerate}

\subsection*{Packet 4: Fair Automated Decisions (Questions 22--27)}

Public agencies increasingly use automated systems to rank applications for permits,
benefits, and inspections. One influential account, outcome parity, calls a system fair when
approval and error rates are equal across specified groups. This account usefully directs
attention to patterns that case-by-case review may miss. Yet the chosen groups and time
period can change the measured result, and equal rates do not reveal whether any particular
decision rests on accurate information. Outcome parity therefore supplies a diagnostic,
not a complete account of justified decision making.

A second view, rule disclosure, locates fairness in public access to the code, data fields,
and ranking criteria. It correctly insists that official decisions should be constrained by
reasons that can be stated. But disclosure focuses on what an outside observer can inspect.
An applicant may learn exactly how an erroneous address reduced a score and still have no
recognized way to correct the address. The theory describes visible machinery while leaving
the affected person in a passive position.

A better account begins by distinguishing two senses of transparency. A process may be
transparent because its mechanics are inspectable, or because its decisions are answerable
to those subject to them. The latter requires more than publication. Under a contestability
account, an agency must give a case-specific explanation, provide a forum in which the person
can challenge relevant data or reasoning, and issue a corrected decision when the challenge
succeeds. Here a ``forum'' need not be a courtroom; it is any established procedure with
authority to hear and remedy a challenge.

Contestability does not make group statistics or published rules unimportant. Statistical
patterns can identify where challenges deserve special attention, and disclosed rules make
objections more informed. But those devices support fairness only when joined to an
institutional duty to answer. The central question is not simply whether observers can see
the system or whether aggregate numbers match. It is whether a person governed by a decision
can require the agency to justify and, when necessary, revise it.

\begin{enumerate}\setcounter{enumi}{21}
% Source map: 2024 Feb RC Q22 -- main point.
\item Which choice best states the main point of the passage?
\begin{enumerate}
\item Outcome parity is the only reliable measure of automated fairness.
\item Publishing source code automatically makes every decision fair.
\item Contestability offers a fuller account of fair automated decisions than parity or disclosure alone.
\item Automated systems should never be used by public agencies.
\item Courts are the only legitimate forums for correcting data.
\end{enumerate}

% Source map: 2024 Feb RC Q23 -- criticism of a theory.
\item According to the author, rule disclosure leaves unexplained
\begin{enumerate}
\item why programmers use data fields
\item how published criteria can be read by observers
\item how an affected person can obtain correction of an erroneous decision
\item why agencies compare approval rates
\item whether source code can be stored electronically
\end{enumerate}

% Source map: 2024 Feb RC Q24 -- implication.
\item Which statement is most strongly implied by the outcome-parity account as described?
\begin{enumerate}
\item A system's fairness assessment can vary with the groups and period selected for comparison.
\item Every equal approval rate proves that all individual data are accurate.
\item A fair system must disclose its complete source code.
\item Applicants must be permitted to challenge every ranking criterion.
\item Group statistics have no role in identifying possible problems.
\end{enumerate}

% Source map: 2024 Feb RC Q25 -- purpose of distinction.
\item The author distinguishes inspectable transparency from answerable transparency primarily to
\begin{enumerate}
\item show why group statistics are mathematically invalid
\item separate mere visibility of rules from institutional responsibility to affected people
\item argue that published code is always misleading
\item establish that applicants should design agency software
\item prove that only private decisions can be transparent
\end{enumerate}

% Source map: 2024 Feb RC Q26 -- meaning in context.
\item As used in the third paragraph, ``forum'' most nearly means
\begin{enumerate}
\item an authorized procedure for presenting and remedying a challenge
\item a public website containing general discussion
\item a schedule for software maintenance
\item the physical office where programmers work
\item an election among affected applicants
\end{enumerate}

% Source map: 2024 Feb RC Q27 -- author evaluation.
\item The author would most likely say that rule disclosure
\begin{enumerate}
\item identifies the importance of stated reasons but mistakes inspectability for sufficient accountability
\item is worthless because published rules can never inform a challenge
\item fully explains how agencies should correct individual errors
\item should replace both statistics and contestability
\item proves that every automated system treats groups equally
\end{enumerate}
\end{enumerate}

\newpage
\section*{Answer Key and Concise Explanations}

\subsection*{Packet 1}
\begin{enumerate}
\item \textbf{C.} The author uses results and design philosophy to present adaptive restoration as a model.
\item \textbf{B.} Vega explicitly rejects restoration as reproduction of one historical snapshot.
\item \textbf{A.} The passage assigns oxygen introduction to the gravel riffles.
\item \textbf{D.} The final paragraph identifies measurable gains plus revision capacity as the greatest achievement.
\item \textbf{E.} Fish were monitored, but the passage never says that hatchery fish were added.
\item \textbf{C.} Vega treats observation-driven revision as integral rather than as evidence of failure.
\end{enumerate}

\subsection*{Packet 2}
\begin{enumerate}\setcounter{enumi}{6}
\item \textbf{B.} Both programs combine access to instruments with help interpreting the resulting data.
\item \textbf{D.} Only Passage B describes a mobile laboratory following a regular route.
\item \textbf{A.} The phrase marks the narrower starting model before calibration, context, and review were added.
\item \textbf{B.} Prior commitment supplies a competing explanation for unusually high retention.
\item \textbf{A.} Each passage argues that equipment alone is insufficient for consistently useful community evidence.
\item \textbf{D.} Calibration, field context, dashboards, and rapid feedback all make readings interpretable.
\item \textbf{C.} Passage A presents lending as valuable but substantially improved by added support.
\item \textbf{C.} The later services preserve the useful core while adding complementary guidance.
\end{enumerate}

\subsection*{Packet 3}
\begin{enumerate}\setcounter{enumi}{14}
\item \textbf{B.} The passage weighs two live explanations, giving evidence and a difficulty for each.
\item \textbf{B.} The third paragraph states that related microbes create the dark film when oxygen is scarce.
\item \textbf{A.} The paragraph introduces the hearth account, its supporting dates, and its carbon shortfall.
\item \textbf{B.} The last sentence treats local testability as an advantage despite radical implications.
\item \textbf{D.} Both proposals revise a basic model yet generate evidence accessible to present testing.
\item \textbf{A.} The proposed isotope comparison expressly uses young nearby deposits to investigate the ancient mechanism.
\item \textbf{B.} Matching particles plus independent temporal alignment directly connects the bands to ancient combustion.
\end{enumerate}

\subsection*{Packet 4}
\begin{enumerate}\setcounter{enumi}{21}
\item \textbf{C.} The author treats parity and disclosure as useful but argues that answerable, revisable decisions are essential.
\item \textbf{C.} The erroneous-address example shows the missing route from visibility to correction.
\item \textbf{A.} The first paragraph expressly notes sensitivity to the comparison groups and period.
\item \textbf{B.} The distinction prepares the move from observable mechanics to enforceable institutional response.
\item \textbf{A.} The passage directly defines a forum as an authorized process that can hear and remedy a challenge.
\item \textbf{A.} Rule disclosure recognizes a genuine need for stated reasons but leaves affected people passive.
\end{enumerate}

\endgroup
